\begin{description}
  \item[Q1.] \textbf{Does the \(\varepsilon^{2q+1}\) identity hold under a
    biased or non-uniform source preparation \(U\ne H^{\otimes n}\)?}
    Basis: our verification used the canonical Grover source
    \(U = H^{\otimes n}\); the paper states the result in terms of an
    arbitrary \(U\), but doesn't quantify how sensitivity to a
    non-uniform \(U\) affects the residual error's exact match to
    \(\varepsilon^{2q+1}\). Next steps: run the same simulator with
    \(U\) sampled from Haar / QR-decomposed random unitaries at
    \(n\!=\!3,4\), holding \(\varepsilon=|U_{ts}|^2\) fixed, and
    check that the identity remains exact rather than approximate.

  \item[Q2.] \textbf{How much does the average-case query saving over
    Phase-\(\pi/3\) actually amount to, as a function of the
    \(\varepsilon\)-distribution?}
    Basis: the paper claims strictly better average behavior; our
    replication tested worst-case only. Next steps: implement Phase-\(\pi/3\)
    as a matching numpy statevector simulator, then Monte-Carlo the two
    algorithms over (i) uniform \(\varepsilon\in(0,1)\), (ii) Beta(2,2),
    and (iii) power-law priors, and compare \(\mathbb{E}[\text{queries to
    reach residual error }10^{-3}]\).

  \item[Q3.] \textbf{Under a realistic depolarizing / amplitude-damping
    noise channel on the ancillas, at what noise level does the
    fixed-point property survive?}
    Basis: the paper explicitly cites "any errors due to imperfect
    transformations in earlier iterations are wiped out by subsequent
    iterations" as a fixed-point feature; but this is stated qualitatively
    only. Next steps: convert our simulator to a density-matrix simulator,
    insert per-gate depolarizing noise with parameter \(p\), and measure
    the ratio of residual error to the ideal \(\varepsilon^{2q+1}\) as a
    function of \(p\) and \(q\).

  \item[Q4.] \textbf{Is the two-ancilla design tight, or can a single-ancilla
    +mid-circuit-measurement variant achieve \(\varepsilon^{2q+1}\) for all
    \(q\) as well?}
    Basis: the paper motivates the second ancilla purely as a device to
    cap \(f\) at 1/2 to enter the simple-scheme's dominance regime; it
    does not prove that one ancilla plus adaptive measurement cannot
    also achieve the all-\(q\) optimum. Next steps: exhaustive search
    over 1-ancilla + adaptive-measurement circuit families up to \(q=3\)
    at \(n=2\), scoring against \(\varepsilon^{2q+1}\).

  \item[Q5.] \textbf{What is the (empirically-observed) rounding-error growth
    in \(\varepsilon^{2q+1}\) vs measured error as a function of $q$, and
    does it match the theoretical stability bound?}
    Basis: our data show a monotone rise in \(|\Delta|\) from
    \(\sim 3\!\times\!10^{-16}\) at \(q=1\) to \(\sim 3\!\times\!10^{-15}\)
    at \(q=4\), suggesting linear-in-\(q\) accumulation of unit-roundoff
    error rather than the multiplicative growth one might naively expect
    from nested renormalisations. Next steps: push \(q\) to 20 at \(n=2\)
    (still cheap), fit \(|\Delta|\) vs \(q\), and compare against the
    condition-number bound for the joint-diffusion reflection.
\end{description}
